%%=============================================================================
%% Conclusie
%%=============================================================================

\chapter{Conclusie}%
\label{ch:conclusie}

% TODO: Trek een duidelijke conclusie, in de vorm van een antwoord op de
% onderzoeksvra(a)g(en). Wat was jouw bijdrage aan het onderzoeksdomein en
% hoe biedt dit meerwaarde aan het vakgebied/doelgroep? 
% Reflecteer kritisch over het resultaat. In Engelse teksten wordt deze sectie
% ``Discussion'' genoemd. Had je deze uitkomst verwacht? Zijn er zaken die nog
% niet duidelijk zijn?
% Heeft het onderzoek geleid tot nieuwe vragen die uitnodigen tot verder 
%onderzoek?

Dit onderzoek vertrok vanuit de vraag hoe de HOGENT-campusapp inclusiever kan worden ontworpen voor studenten met een beperking, en welke impact dergelijke aanpassingen hebben op de performantie van de applicatie. Op basis van zowel geautomatiseerde audits als manuele evaluaties kan geconcludeerd worden dat gerichte toegankelijkheidsverbeteringen op UI- en componentniveau leiden tot een aantoonbaar betere bruikbaarheid, zonder significante negatieve impact op de performantie.

De resultaten tonen aan dat de oorspronkelijke applicatie, ondanks een visueel verzorgde opbouw, belangrijke tekortkomingen vertoonde op het vlak van toegankelijkheid. Zoals vaak het geval is, betekent een verzorgd ontwerp niet automatisch dat een applicatie ook werkelijk toegankelijk is voor alle gebruikers. In het bijzonder werden problemen vastgesteld met semantische structuur, toetsenbordnavigatie en ondersteuning voor schermlezers. Door deze aspecten systematisch te analyseren en te verbeteren — onder andere via een correcte focusvolgorde, semantische HTML en consistente interactiepatronen — werd de applicatie aanzienlijk toegankelijker gemaakt. Dit resulteerde in een interface die beter bruikbaar is voor gebruikers met een visuele of motorische beperking, en tegelijk ook de algemene gebruikservaring verbeterde.

De vergelijking tussen nulmeting en hertesting bevestigt deze evolutie. Zowel de Lighthouse-metingen als de manuele evaluaties tonen aan dat de applicatie beter voldoet aan de WCAG-richtlijnen en correcter wordt ondersteund door assistieve technologieën zoals schermlezers. Tegelijk blijkt uit de performantiemetingen dat de impact van deze aanpassingen beperkt blijft. Hoewel individuele metrics kunnen variëren afhankelijk van testomstandigheden en meetmomenten, wijzen de resultaten niet op structurele performantieproblemen. Dit ondersteunt de vaststelling dat toegankelijkheid en performantie geen tegenstrijdige doelstellingen zijn, maar samen kunnen bijdragen aan de kwaliteit van een applicatie.

De belangrijkste bijdrage van dit onderzoek ligt in het aantonen dat toegankelijkheid niet als een bijkomende of optionele factor mag worden beschouwd, maar als een fundamenteel onderdeel van softwarekwaliteit. Door systematisch te testen, problemen te identificeren, verbeteringen te implementeren en deze opnieuw te evalueren, werd een iteratieve werkwijze gevolgd die aansluit bij moderne ontwikkelpraktijken. Dit is in het bijzonder relevant binnen een agile context, waar toegankelijkheidsverbeteringen geïntegreerd kunnen worden in opeenvolgende ontwikkelcycli in plaats van pas op het einde te worden toegevoegd.

Kritisch beschouwd blijven er nog enkele aandachtspunten. Performantiemetingen blijven gevoelig voor externe factoren en zouden in toekomstig onderzoek verder verfijnd kunnen worden via herhaalde en gecontroleerde testopstellingen. Daarnaast kan de toegankelijkheid nog uitgebreider worden onderzocht door bijkomende gebruikersgroepen en assistieve technologieën te betrekken, evenals door meer diepgaande gebruikerstesten in realistische contexten.

Samenvattend toont deze bachelorproef aan dat digitale toegankelijkheid een fundamentele vereiste is binnen moderne applicatieontwikkeling. De uitgevoerde optimalisaties leveren een duidelijke meerwaarde op voor de eindgebruiker, in het bijzonder voor studenten met een functiebeperking, terwijl de technische impact beperkt blijft. De praktische meerwaarde van dit onderzoek ligt dus niet enkel in de verbeterde applicatie zelf, maar ook in de vaststelling dat toegankelijkheidsbewust ontwikkelen een haalbare en noodzakelijke praktijk is binnen moderne softwareprojecten.

% Dit onderzoek vertrok vanuit de vraag hoe de digitale toegankelijkheid van de betrokken applicatie kon worden verbeterd, en welke impact die verbeteringen zouden hebben op de algemene kwaliteit en performantie van het systeem. De resultaten tonen duidelijk aan dat toegankelijkheid niet los gezien kan worden van softwarekwaliteit: wanneer semantische structuur, contrast, navigatie en gebruiksvriendelijkheid doelgericht worden aangepakt, stijgt de bruikbaarheid van de applicatie merkbaar voor een bredere groep gebruikers. De combinatie van automatische Lighthouse-metingen en manuele controles maakte het mogelijk om niet alleen vast te stellen dat de toegankelijkheid verbeterde, maar ook om te begrijpen waarom die verbeteringen relevant zijn in de praktijk.

% De belangrijkste conclusie is dat de uitgevoerde aanpassingen een duidelijke meerwaarde opleveren voor de eindgebruiker. De applicatie was al vrij goed opgebouwd vanuit een designperspectief, maar zoals vaak het geval is, betekent een verzorgd ontwerp niet automatisch dat een applicatie ook werkelijk toegankelijk is voor iedereen. Door gericht te testen aan de hand van WCAG-principes en Lighthouse-audits werd zichtbaar waar de pijnpunten zaten en hoe die konden worden weggewerkt. Vooral de verbeteringen rond semantische HTML, kleurcontrast, structurele opbouw en consistente interactiepatronen hebben ervoor gezorgd dat de applicatie beter bruikbaar is voor mensen met visuele of motorische beperkingen, en in bredere zin ook voor gebruikers die voordeel halen uit een duidelijkere en voorspelbare interface.

% De performantie-impact van die toegankelijkheidsaanpassingen bleef beperkt. Dat is een belangrijk resultaat, omdat het aantoont dat toegankelijkheid en performantie geen tegenstrijdige doelen hoeven te zijn. In deze case bleef de visuele en functionele kwaliteit van de applicatie overeind, terwijl de toegankelijkheid aanzienlijk vooruitging. De metingen laten wel zien dat performantiemetrics kunnen variëren afhankelijk van de testcontext, het meetmoment en mogelijke uitschieters. Daarom is het belangrijk om performantie niet uitsluitend op één meetresultaat te beoordelen, maar eerder te kijken naar patronen over meerdere tests heen. Dat nuanceert de interpretatie van de cijfers, maar verandert niets aan de centrale vaststelling dat de toegankelijkheidswinst reëel en waardevol is.

% Mijn bijdrage aan dit onderzoek bestond er vooral in om toegankelijkheid niet als een bijkomend detail te behandelen, maar als een essentieel kwaliteitsaspect van de applicatie. Door systematisch te testen, problemen te identificeren, verbeteringen te implementeren en die opnieuw te evalueren, werd een iteratieve werkwijze gevolgd die goed aansluit bij hoe software in de praktijk ontwikkeld wordt. Die aanpak heeft ook een bredere meerwaarde voor het vakgebied en voor de doelgroep: ze toont aan dat toegankelijkheid niet noodzakelijk aan het einde van een project moet worden bekeken, maar dat ze vanaf het begin en doorheen het hele ontwikkelproces geïntegreerd kan worden.

% Dat is vooral relevant binnen een agile context. Toegankelijkheidsverbeteringen kunnen perfect worden meegenomen in iteraties, net zoals andere functionele of technische verbeteringen. Door toegankelijkheid in elke sprint of ontwikkelfase mee te nemen, vermijd je dat het onderwerp pas op het einde wordt gecontroleerd of als een last-minute correctie moet worden opgelost. Dat maakt het proces efficiënter en vergroot de kans dat toegankelijkheid structureel ingebed raakt in de manier waarop teams software bouwen. In die zin is dit onderzoek niet alleen een evaluatie van één applicatie, maar ook een illustratie van hoe ontwikkelteams toegankelijkheid als standaardonderdeel van hun werkwijze kunnen opnemen.

% De uitkomst van dit onderzoek was deels verwacht, maar niet volledig vanzelfsprekend. Het was vooraf aannemelijk dat de applicatie al een degelijke basis zou hebben, zeker omdat het ontwerp en de structuur al relatief verzorgd waren. Tegelijk toonde de evaluatie aan dat zelfs een goed ogende applicatie nog belangrijke toegankelijkheidsverbeteringen nodig kan hebben. Dat bevestigt een belangrijk inzicht: visuele kwaliteit is niet hetzelfde als toegankelijkheid. Zonder specifieke testmethodes zoals Lighthouse, WCAG-controles en manuele verificatie blijven veel problemen onopgemerkt. Net die combinatie van automatische en manuele analyse bleek noodzakelijk om een correct beeld te krijgen van de werkelijke kwaliteit van de applicatie.

% Kritisch bekeken blijven er uiteraard nog aandachtspunten. Ten eerste kunnen performantiemetingen altijd verder verfijnd worden door herhaalde tests onder streng gecontroleerde omstandigheden. Ten tweede kan toegankelijkheid nog breder onderzocht worden door extra gebruikersgroepen en ondersteunende technologieën mee te nemen, bijvoorbeeld schermlezers, toetsenbordnavigatie of situaties met cognitieve beperkingen. Ook op vlak van inhoud, interactie en responsief gedrag zijn er nog mogelijkheden om de applicatie verder te optimaliseren. Het onderzoek heeft dus niet enkel antwoorden opgeleverd, maar ook nieuwe vragen geopend over hoe toegankelijkheid nog verder kan worden uitgebreid en verfijnd.

% Daaruit volgt een duidelijke eindconclusie: toegankelijkheid moet beschouwd worden als een standaardvereiste binnen applicatieontwikkeling, niet als een optionele extra. De resultaten van deze bachelorproef tonen aan dat toegankelijkheidsverbeteringen substantieel verschil maken voor de gebruikservaring, terwijl de impact op performantie beperkt blijft. Bovendien is aangetoond dat dergelijke verbeteringen haalbaar zijn binnen een iteratief ontwikkelproces en dus perfect geïntegreerd kunnen worden in agile softwareontwikkeling. Voor ontwikkelteams betekent dit dat toegankelijkheid vanaf het begin mee moet worden opgenomen, zodat digitale producten niet alleen functioneel en aantrekkelijk zijn, maar ook werkelijk bruikbaar voor zoveel mogelijk mensen.

% Samengevat laat dit onderzoek zien dat een applicatie pas echt kwalitatief is wanneer ze technisch goed werkt, toegankelijk is voor diverse gebruikers en op een verantwoorde manier verder kan worden verbeterd. De praktische meerwaarde ligt dus niet alleen in de aangepaste applicatie zelf, maar ook in de vaststelling dat toegankelijkheidsbewust ontwikkelen een haalbare en noodzakelijke gewoonte is voor moderne softwareprojecten.



